\section{投资框架：如何看 Coding/Agent 公司的价值}

本期最后谈到投资新思考。若 Coding 是 AI 加速器，投资判断就不能只看用户增长或模型榜单，而要看一家公司是否掌握工作流、反馈和成本结构。一个 Coding/Agent 公司如果只是套壳调用模型，壁垒很薄；如果能积累真实任务数据、失败案例、工具集成和用户工作流，壁垒会更深。

\subsection{四个判断维度}

\noindent\begin{tabularx}{\textwidth}{p{0.23\textwidth}p{0.34\textwidth}X}
\toprule
维度 & 要问的问题 & 好信号 \\
\midrule
工作流深度 & 是否进入用户每天必须完成的任务？ & 代码库、测试、部署、数据流程深度接入。 \\
反馈闭环 & 是否积累可训练的失败和成功轨迹？ & 有结构化日志、评估和人类修正数据。 \\
成本结构 & 是否能用多模型路由降低单位任务成本？ & 大模型只做关键步骤，小模型处理高频任务。 \\
分发入口 & 是否拥有用户习惯或生态位？ & IDE、企业系统、浏览器、团队协作入口。 \\
\bottomrule
\end{tabularx}

\begin{warningbox}{投资叙事里的常见陷阱}
“我们是 Agent 公司”不是壁垒；“我们接了最新模型”也不是壁垒。真正的壁垒通常在工作流深度、数据反馈、组织执行和分发入口，而不是 prompt 包装。
\end{warningbox}

\subsection{本章小结}

Coding/Agent 投资要从 demo 转向 workflow。能否沉淀反馈、降低成本、占住入口，比短期模型能力展示更重要。

